\chapter{Mesothelioma}
\index{Mesothelioma}

\section*{Definition and Overview}
Mesothelioma is an aggressive, highly lethal malignancy arising from the mesothelial cells that line the body's serous cavities, most commonly the pleura (pleural mesothelioma, representing over 80--90\% of cases) and less frequently the peritoneum (peritoneal mesothelioma), pericardium, or tunica vaginalis. The disease is almost exclusively caused by exposure to asbestos fibres, which migrate to the serosal membranes and induce chronic inflammation, cellular damage, and neoplastic transformation. Mesothelioma is characterised by an exceptionally long latency period, often manifesting 20 to 50 years after the initial asbestos inhalation. Due to its infiltrative nature and frequent resistance to standard oncological therapies, the disease carries a poor prognosis and requires a multidisciplinary, specialised approach to management.

\section*{Epidemiology}
Mesothelioma is a major public health concern, particularly in nations with a history of intensive asbestos mining, manufacturing, or industrial application. some countries have one of the highest age-standardised rates of mesothelioma globally, a consequence of the extensive use of crocidolite (blue asbestos) and chrysotile (white asbestos) in building and industrial materials during the mid-to-late 20th century. Approximately 700 to 800 new cases are diagnosed in many countries each year. The disease shows a marked gender bias, with males representing over 80\% of cases due to historical occupational exposures in trades such as construction, shipbuilding, railway work, and mining. The median age at diagnosis is approximately 75 years, reflecting the long latency period.

\section*{Aetiology and Pathophysiology}
The primary driver of mesothelioma is the physical and chemical interaction of asbestos fibres with mesothelial tissues:
\begin{itemize}
    \item \textbf{Asbestos inhalation:} Microscopic, durable asbestos fibres are inhaled into the alveoli, penetrate the pulmonary parenchyma, and migrate via lymphatic pathways to the pleural space.
    \item \textbf{Chronic inflammation:} The physical presence of long, needle-like asbestos fibres causes mechanical irritation and prevents phagocytosis by macrophages (``frustrated phagocytosis''), leading to the chronic release of mutagenic reactive oxygen species (ROS), cytokines, and growth factors.
    \item \textbf{Genetic alterations:} Chromosomal damage and deletion of key tumour suppressor genes, including \textit{BAP1} (BRCA1-associated protein 1), \textit{NF2} (neurofibromatosis type 2), and \textit{CDKN2A} (cyclin-dependent kinase inhibitor 2A).
    \item \textbf{Germline predisposition:} Inherited mutations in the \textit{BAP1} gene increase susceptibility to developing mesothelioma, even with low levels of environmental asbestos exposure.
    \item \textbf{Non-asbestos mineral fibres:} Inhalation of other fibrous silicate minerals such as erionite (prevalent in parts of Turkey) or fluoro-edenite can also induce mesothelioma.
\end{itemize}

\section*{Clinical Features}
The clinical features of mesothelioma are insidious and progressive, typically worsening over several months prior to diagnosis:
\begin{itemize}
    \item \textbf{Progressive dyspnoea:} Shortness of breath, initially on exertion and later at rest, caused by a large, recurrent pleural effusion or restriction of lung expansion by a thick, encasing tumour sheet.
    \item \textbf{Chest wall pain:} Persistent, localized, dull, or aching chest pain, which may become pleuritic or radiate to the shoulder, due to invasion of the intercostal nerves and ribs.
    \item \textbf{Pleural effusion signs:} Decreased breath sounds, dullness to percussion, and reduced vocal fremitus on the affected side of the chest.
    \item \textbf{Constitutional symptoms:} Unintentional and progressive weight loss, anorexia, fatigue, low-grade fevers, and night sweats.
    \item \textbf{Abdominal symptoms (peritoneal):} Abdominal pain, progressive abdominal distension (due to ascites), early satiety, nausea, and occasionally palpable abdominal masses or bowel obstruction.
    \item \textbf{Signs of local invasion:} Dysphagia (due to oesophageal compression), hoarseness (due to recurrent laryngeal nerve involvement), or superior vena cava syndrome (facial and upper limb swelling).
\end{itemize}

\section*{Differential Diagnosis}
Because the clinical and radiological features of mesothelioma mimic other pulmonary and pleural conditions, the differential diagnoses include:
\begin{itemize}
    \item \textbf{Primary lung adenocarcinoma:} Particularly cases presenting with pleural metastases, distinguished by an immunohistochemical (IHC) panel (mesothelioma is positive for calretinin, WT1, and D2-40, and negative for TTF-1 and CEA).
    \item \textbf{Metastatic pleural carcinoma:} Metastases from breast, ovarian, or renal carcinomas, differentiated by cytological and histopathological immunophenotyping.
    \item \textbf{Benign asbestos pleural effusion (BAPE):} A self-limiting, non-malignant, exudative effusion occurring in individuals with prior asbestos exposure, distinguished by the absence of atypical cells on cytology and resolution without recurrence.
    \item \textbf{Pleural empyema or tuberculous pleurisy:} Infectious processes causing pleural thickening and fluid accumulation, distinguished by bacterial cultures, PCR, and pleural biopsy.
    \item \textbf{Asbestos-related pleural plaques:} Benign, localized, calcified areas of pleural thickening that do not invade the lung parenchyma or chest wall and are asymptomatic.
\end{itemize}

\section*{When to Seek Medical Care}
Individuals with a known history of occupational or environmental asbestos exposure should seek prompt medical advice if they develop persistent chest pain, a new cough, or progressive shortness of breath.

\textbf{Contact your local emergency services for:}
\begin{itemize}
    \item Sudden, severe difficulty breathing or a feeling of suffocation.
    \item Sharp, crushing chest pain or chest tightness.
    \item Coughing up blood (haemoptysis).
    \item Signs of severe respiratory distress, such as blue lips (cyanosis), confusion, or sudden collapse.
\end{itemize}

\section*{Diagnosis and Investigations}
Establishing a diagnosis of mesothelioma requires advanced imaging, tissue biopsies, and detailed immunohistochemical staining:
\begin{itemize}
    \item \textbf{Chest X-ray:} Often the initial investigation, showing a unilateral pleural effusion, pleural thickening, or volume loss of the affected hemithorax.
    \item \textbf{Contrast-enhanced chest and abdominal CT scan:} Demonstrates nodular, circumferential pleural thickening, pleural plaques, and evaluates local invasion into the chest wall, mediastinum, or diaphragm.
    \item \textbf{VATS pleural biopsy (Thoracoscopy):} The gold standard diagnostic method, allowing direct visualization of the pleura, drainage of effusion fluid, and acquisition of multiple deep tissue biopsies containing both pleural and underlying lung tissue.
    \item \textbf{Pleural fluid cytology:} Thoracocentesis allows cytological analysis, but this has a low diagnostic sensitivity (30--50\%) and is rarely sufficient for a definitive diagnosis of mesothelioma.
    \item \textbf{Immunohistochemistry (IHC) panel:} Essential to confirm mesothelial origin (using positive markers calretinin, Wilms tumour protein-1 [WT1], cytokeratin 5/6, and podoplanin [D2-40]) and to exclude adenocarcinoma.
    \item \textbf{PET-CT scan:} Used to assess metabolic activity, stage the disease, guide biopsy site selection, and monitor response to oncology treatments.
\end{itemize}

\section*{Management}
Management is palliative and focuses on extending survival and controlling symptoms using a combination of systemic therapies and local interventions:
\begin{itemize}
    \item \textbf{Systemic immunotherapy:} First-line combination immunotherapy (nivolumab plus ipilimumab) is now standard for unresectable mesothelioma, providing significant survival benefits, particularly in non-epithelioid subtypes.
    \item \textbf{Systemic chemotherapy:} Platinum-based chemotherapy (cisplatin or carboplatin combined with pemetrexed) remains a primary treatment option.
    \item \textbf{Pleural fluid management:} Recurrent symptomatic pleural effusions are managed with chemical pleurodesis (using sterile talc slurry to adhere the pleural layers) or the insertion of an indwelling pleural catheter (IPC) for home drainage.
    \item \textbf{Palliative radiotherapy:} Directed to localized areas of tumour growth to manage severe chest wall pain or to prevent tumour tract seeding after biopsy.
    \item \textbf{Surgical cytoreduction:} Highly selected patients may undergo pleurectomy and decortication (P/D) to remove gross tumour bulk, though extrapleural pneumonectomy (EPP) is now rarely performed due to high morbidity.
    \item \textbf{Early palliative care integration:} Early referral to palliative specialists to manage neuropathic pain, dyspnoea, and support quality of life.
\end{itemize}

\section*{Complications}
The progression of mesothelioma leads to severe, localized thoracic and systemic complications:
\begin{itemize}
    \item \textbf{Severe respiratory failure:} Caused by the encasement of the lung (``frozen chest'') and progressive loss of functional lung volume.
    \item \textbf{Superior Vena Cava (SVC) obstruction:} Compression of the superior vena cava by mediastinal tumour invasion, causing venous congestion and swelling of the head and arms.
    \item \textbf{Intractable neuropathic pain:} Resulting from direct tumor invasion of the intercostal nerves, brachial plexus, or chest wall bones.
    \item \textbf{Tumour cachexia:} Severe wasting syndrome with muscle loss and anorexia driven by chronic inflammatory cytokines.
    \item \textbf{Pericardial tamponade:} Infiltration of the pericardial space by the tumour, restricting cardiac filling and causing cardiovascular collapse.
\end{itemize}

\section*{Prognosis and Outcomes}
The prognosis for mesothelioma is poor. The disease is highly aggressive, and the median survival time from diagnosis is typically 9 to 18 months. The overall 5-year survival rate is low, ranging between 5\% and 10\%. Prognosis is heavily influenced by the histological subtype: epithelioid mesothelioma has a significantly better prognosis and response to treatment than sarcomatoid or biphasic subtypes. Other favourable prognostic factors include early-stage disease, good performance status, younger age, and eligibility for dual immunotherapy.

\section*{Prevention and Self-Care}
Preventing mesothelioma depends entirely on avoiding exposure to asbestos fibres:
\begin{itemize}
    \item \textbf{Workplace safety compliance:} Adhere strictly to occupational health and safety regulations if working in construction, plumbing, electrical, or demolition trades where asbestos may be present.
    \item \textbf{Licensed asbestos removal:} Never disturb, drill, cut, or attempt to remove asbestos-containing materials (such as fibro sheets) in domestic properties; always engage a licensed asbestos removalist.
    \item \textbf{Personal protective equipment (PPE):} Use certified particulate respirators (HEPA-filtered) and disposable overalls when working in environments with potential mineral dust.
    \item \textbf{Smoking cessation:} While tobacco smoke does not cause mesothelioma, it works synergistically with asbestos to dramatically increase the risk of bronchogenic lung cancer and worsens baseline lung function.
    \item \textbf{Preventative health care:} Stay up to date with seasonal influenza and pneumococcal vaccinations to reduce the risk of secondary respiratory infections.
\end{itemize}

\section*{Sources and Further Reading}
The following reputable organisations provide further information on mesothelioma:
\begin{itemize}
    \item Cancer Council Australia. \textit{Information on mesothelioma, financial support, and support groups.} https://www.cancer.org.au/
    \item Asbestos Diseases Research Institute (ADRI). \textit{Clinical research updates and patient support for asbestos-related illnesses.} https://adri.org.au/
    \item NHS. \textit{Overview of mesothelioma diagnosis, symptoms, and treatment.} https://www.nhs.uk/
    \item Mayo Clinic. \textit{Details of mesothelioma staging, surgical options, and clinical trials.} https://www.mayoclinic.org/
    \item National Cancer Institute (NCI). \textit{Comprehensive guide to pleural and peritoneal mesothelioma.} https://www.cancer.gov/
\end{itemize}
